\documentclass[a4paper]{report}

\usepackage[
    theme = epfl-notes-legacy,
    palette = dark
]{latex-fayorg-template}

\ProjectSetup{
    title = Notes Example,
    subtitle = {A showcase of reusable note components},
    author = Elie Baier (397222),
    context = Computer Science Bachelor,
    course_code = CS-101,
    course_name = {Advanced Information, Computation and Communication II},
    course_semester = Spring 2026,
    course_instructor = Christoph Gastpar
}

\begin{document}
    \MakeProjectFrontPage

    \section{Introduction}
    This is a simple example of a LaTeX document using \ProjectTitle (for \ProjectCourseCode, \ProjectCourseName) created by \ProjectAuthor.

    \begin{theorem}[A useful theorem]
        If $a=b$ and $b=c$, then $a=c$.
    \end{theorem}

    \begin{definition}[A reusable component]
        A semantic component keeps its meaning while its appearance follows the active theme.
    \end{definition}

    \begin{example}
        Select the personal theme without changing the contents of this environment.
    \end{example}

    \begin{remark}
        Remarks are unnumbered in the legacy setup.
    \end{remark}

    \begin{exercise}[Short][Transitivity]
        Prove the theorem above.
        \Answer
        Apply equality twice.
    \end{exercise}

    \begin{explanation}
        This unnumbered proof-style box ends with the legacy marker.
    \end{explanation}

    \begin{conjecture}
        The legacy component API can still evolve without changing old notes.
    \end{conjecture}

    \begin{answer}
        Standalone answer boxes are also available.
    \end{answer}

    \begin{note}
        Notes are unnumbered and intended for short callouts.
    \end{note}

    \begin{code}[language=Python]
def short_commit(commit):
    return commit[:7]
    \end{code}
\end{document}
